% Options for packages loaded elsewhere
\PassOptionsToPackage{unicode}{hyperref}
\PassOptionsToPackage{hyphens}{url}
%
\documentclass[
]{article}
\usepackage{amsmath,amssymb}
\usepackage{iftex}
\ifPDFTeX
  \usepackage[T1]{fontenc}
  \usepackage[utf8]{inputenc}
  \usepackage{textcomp} % provide euro and other symbols
\else % if luatex or xetex
  \usepackage{unicode-math} % this also loads fontspec
  \defaultfontfeatures{Scale=MatchLowercase}
  \defaultfontfeatures[\rmfamily]{Ligatures=TeX,Scale=1}
\fi
\usepackage{lmodern}
\ifPDFTeX\else
  % xetex/luatex font selection
\fi
% Use upquote if available, for straight quotes in verbatim environments
\IfFileExists{upquote.sty}{\usepackage{upquote}}{}
\IfFileExists{microtype.sty}{% use microtype if available
  \usepackage[]{microtype}
  \UseMicrotypeSet[protrusion]{basicmath} % disable protrusion for tt fonts
}{}
\makeatletter
\@ifundefined{KOMAClassName}{% if non-KOMA class
  \IfFileExists{parskip.sty}{%
    \usepackage{parskip}
  }{% else
    \setlength{\parindent}{0pt}
    \setlength{\parskip}{6pt plus 2pt minus 1pt}}
}{% if KOMA class
  \KOMAoptions{parskip=half}}
\makeatother
\usepackage{xcolor}
\usepackage{longtable,booktabs,array}
\usepackage{calc} % for calculating minipage widths
% Correct order of tables after \paragraph or \subparagraph
\usepackage{etoolbox}
\makeatletter
\patchcmd\longtable{\par}{\if@noskipsec\mbox{}\fi\par}{}{}
\makeatother
% Allow footnotes in longtable head/foot
\IfFileExists{footnotehyper.sty}{\usepackage{footnotehyper}}{\usepackage{footnote}}
\makesavenoteenv{longtable}
\setlength{\emergencystretch}{3em} % prevent overfull lines
\providecommand{\tightlist}{%
  \setlength{\itemsep}{0pt}\setlength{\parskip}{0pt}}
\setcounter{secnumdepth}{-\maxdimen} % remove section numbering
\ifLuaTeX
  \usepackage{selnolig}  % disable illegal ligatures
\fi
\usepackage{bookmark}
\IfFileExists{xurl.sty}{\usepackage{xurl}}{} % add URL line breaks if available
\urlstyle{same}
\hypersetup{
  pdftitle={Affine Maps Compatible with Modular Rounding},
  pdfauthor={Samir Bouftass and ChatGPT (OpenAI)},
  hidelinks,
  pdfcreator={LaTeX via pandoc}}

\title{Affine Maps Compatible with Modular Rounding}
\usepackage{etoolbox}
\makeatletter
\providecommand{\subtitle}[1]{% add subtitle to \maketitle
  \apptocmd{\@title}{\par {\large #1 \par}}{}{}
}
\makeatother
\subtitle{Exact factorization, missing carries, and prediction loss}
\author{Samir Bouftass and ChatGPT (OpenAI)}
\date{15 November 2026 · English edition 1.1}

\begin{document}
\maketitle

\subsection{Abstract}\label{abstract}

Let \(P=KQ\) with \(K,Q\ge2\), and let \(F(z)=\lfloor[z]_P/Q\rfloor\)
act coordinatewise. We classify all affine maps for which \(F(Uz+b)\)
can be computed from \(F(z)\) for every input. Each row of \(U\),
reduced modulo \(P\), must be zero or contain exactly one nonzero entry,
equal to \(1\) or \(-1\). Positive and negative rows require different
congruences on the translation. In particular, an invertible square
linear map compatible with this rounding is a permutation matrix;
allowing an appropriate translation enlarges the class to signed
permutations. We also quantify the information lost in a two-coordinate
sum. Its optimal prediction error from the rounded inputs is
\((Q-1)/(2Q)\). Resampling the missing carry reproduces the public
marginal exactly, but generally fails to preserve its dependence on the
unrounded input or a hidden secret. The contribution is an explicit
normal form for a fixed interval partition and elementary, sharp
examples of the distinction between public simulation and joint
simulation. The general principles of factorization through fibers and
optimal finite prediction are standard.

\subsection{1. Question and setting}\label{question-and-setting}

Rounding occurs in lattice-based cryptography, including Learning with
Rounding and its variants {[}1--5{]}. A linear operation on an unrounded
value need not remain computable after rounding: the operation may
depend on a remainder that has already been discarded. The question
considered here is exact and algebraic. For which public affine maps
does there exist a function \(g\) such that

\[
F(Uz+b)=g(F(z))\qquad\text{for every }z\in\mathbb Z_P^m?
\tag{1}
\]

We write \([x]_P\) for the representative of \(x\) in
\(\{0,\ldots,P-1\}\). Fix integers \(P=KQ\), \(K,Q\ge2\), and positive
dimensions \(m,n\). The map \(F:\mathbb Z_P^m\to\mathbb Z_K^m\) is
applied coordinatewise. The matrix \(U\) has size \(n\times m\), and
\(b\) has length \(n\); both have integer entries interpreted modulo
\(P\). No linearity, efficiency, or other restriction is imposed on
\(g\).

Equation (1) is a universal identity. It differs from an identity valid
only on a special input support, from a probabilistic approximation, and
from a cryptographic reduction that changes dimensions, secrets, noise,
or moduli. We make these distinctions explicit after proving the
classification.

The fibers of \(F\) form a fixed partition into products of consecutive
intervals. Maps preserving a partition are classical objects {[}6{]};
congruence-preserving maps concern a related but different kind of
partition {[}7{]}. Our result computes the affine maps compatible with
this particular interval partition. It does not introduce the general
fiber-factorization criterion.

\subsection{2. The exact carry identity}\label{the-exact-carry-identity}

For each input, write

\[
z=Qy+r+Ph,\qquad y=F(z),\qquad 0\le r_j<Q,
\]

where \(h\) is an integer vector if an arbitrary integer representative
of \(z\) is used. Euclidean division gives

\[
F(Uz+b)=\left[Uy+\left\lfloor\frac{Ur+b}{Q}\right\rfloor\right]_K.
\tag{2}
\]

Indeed, adding a multiple of \(P=KQ\) changes the integer quotient by a
multiple of \(K\). The remaining multiple \(QUy\) can be taken outside
the floor. This proves (2), including negative coefficients, negative
translations, and modular wraparound. Floors and reductions in this
formula act coordinatewise.

The correction term depends on the discarded remainder \(r\).
Consequently, (1) holds if and only if \(F(Ur+b)\) is constant over the
zero cell \(\{0,\ldots,Q-1\}^m\). Necessity follows by taking \(y=0\).
For sufficiency, (2) then gives the same correction for every \(y\).

\subsection{3. Complete affine
classification}\label{complete-affine-classification}

\textbf{Lemma 1 (a progression contained in one cell).} Suppose
\(a\ne0\) in \(\mathbb Z_P\). If the \(Q\) points \([t+ar]_P\),
\(0\le r<Q\), lie in one aligned interval of length \(Q\), then either
\(a=1\) and \(t\equiv0\pmod Q\), or \(a=-1\) and \(t\equiv Q-1\pmod Q\).

\textbf{Proof.} Set \(d=\gcd(a,P)\) and \(\ell=P/d\), the order of
\(a\). If \(\ell<Q\), the progression contains the entire coset of the
subgroup generated by \(a\). Its representatives have span
\(P-d\ge P/2\ge Q\), so they cannot fit in an interval of span \(Q-1\).
If \(\ell=Q\), the \(Q\) distinct points fill the interval. Since all
have one residue modulo \(d\), two consecutive points force \(d=1\),
contradicting \(\ell=Q<P\).

Thus \(\ell>Q\). The \(Q\) points again fill the interval, forcing
\(d=1\). Write this interval as \(C\). The progressions indexed by
\(0,\ldots,Q-1\) and \(1,\ldots,Q\) have exactly \(Q-1\) common points,
so \(|C\cap(C+a)|=Q-1\). For \(1\le a<P\) and \(Q\le P/2\),

\[
|C\cap(C+a)|=\max(0,Q-a)+\max(0,Q-(P-a)).
\]

At most one term is positive. Equality with \(Q-1\) therefore forces
\(a=1\) or \(a=P-1\). Alignment of the first or last point gives the
stated congruence on \(t\). \(\square\)

\textbf{Theorem 1 (normal form).} The following conditions are
equivalent:

\begin{enumerate}
\def\labelenumi{\arabic{enumi}.}
\tightlist
\item
  There exists \(g\) satisfying (1).
\item
  The map \(r\mapsto F(Ur+b)\) is constant on \(\{0,\ldots,Q-1\}^m\).
\item
  Every row of \(U\), modulo \(P\), has one of the forms in the table
  below, with the indicated condition on \(b_i\).
\end{enumerate}

\begin{longtable}[]{@{}
  >{\raggedright\arraybackslash}p{(\columnwidth - 4\tabcolsep) * \real{0.3333}}
  >{\raggedright\arraybackslash}p{(\columnwidth - 4\tabcolsep) * \real{0.3333}}
  >{\raggedright\arraybackslash}p{(\columnwidth - 4\tabcolsep) * \real{0.3333}}@{}}
\toprule\noalign{}
\begin{minipage}[b]{\linewidth}\raggedright
Row of \(U\)
\end{minipage} & \begin{minipage}[b]{\linewidth}\raggedright
Translation condition
\end{minipage} & \begin{minipage}[b]{\linewidth}\raggedright
Output coordinate \(g_i(y)\) in \(\mathbb Z_K\)
\end{minipage} \\
\midrule\noalign{}
\endhead
\bottomrule\noalign{}
\endlastfoot
Zero & Any \(b_i\) & \(\lfloor[b_i]_P/Q\rfloor\) \\
One entry \(+1\) in column \(j\) & \(b_i\equiv0\pmod Q\) &
\(y_j+\lfloor[b_i]_P/Q\rfloor\) \\
One entry \(-1\) in column \(j\) & \(b_i\equiv Q-1\pmod Q\) &
\(-y_j+\lfloor[b_i]_P/Q\rfloor\) \\
\end{longtable}

Different rows may use the same column. Whenever it exists, \(g\) is
unique.

\textbf{Proof.} The equivalence of the first two conditions follows from
(2). Consider a row satisfying the second condition. Varying one
remainder coordinate while fixing the others at zero and applying Lemma
1 shows that every nonzero coefficient is \(1\) or \(-1\), with the
corresponding congruence on the row translation. A row cannot contain
both signs, because this would require \(b_i\equiv0\) and
\(b_i\equiv Q-1\pmod Q\), which are distinct for \(Q\ge2\).

If two coefficients are \(1\), the translation is \(Q\beta\). Setting
both corresponding remainders to \(Q-1\) changes the correction quotient
from \(\beta\) to \(\beta+1\). If two coefficients are \(-1\), write the
translation as \(Q\beta+Q-1\). The same choice changes the correction
from \(\beta\) to \(\beta-1\). These changes are nonzero modulo
\(K\ge2\), contradicting constancy. Thus at most one coefficient is
nonzero. Substitution into (2) proves sufficiency and the formulas in
the table. Finally, \(F\) is surjective, which gives uniqueness.
\(\square\)

\textbf{Corollary 1.} With \(b=0\), each row is zero or selects one
coordinate with coefficient \(1\). If \(U\) is square and invertible
over \(\mathbb Z_P\), it is a permutation matrix. With unrestricted
admissible translations, square invertible matrices are exactly signed
permutation matrices.

For the last assertion, invertibility rules out zero rows and repeated
selected columns. Conversely, signed permutations are invertible, and
the prescribed translations make their rows admissible. The
classification concerns coefficients modulo \(P\), not merely modulo the
smaller output modulus \(K\).

\subsection{4. Examples and support
restrictions}\label{examples-and-support-restrictions}

\textbf{An invertible shear fails.} Take \(P=4\), \(Q=2\), and the map
\((z_1,z_2)\mapsto(z_1+z_2,z_2)\). Inputs \((0,0)\) and \((1,1)\) have
the same rounded value \((0,0)\). Their transformed rounded values are
\((0,0)\) and \((1,0)\). Thus invertibility before rounding does not
imply computability after rounding.

\textbf{A corrected reflection succeeds.} For every \(z\),

\[
F(-z-1)=K-1-F(z).
\tag{3}
\]

The translation is essential. Unshifted negation does not, in general,
preserve the rounding cells. More generally, translations by multiples
of the quantization step commute with the quotient in the familiar way;
related translation identities also appear in modulo quantization
{[}8{]}.

\textbf{Restricted supports can change the answer.} If every coordinate
of \(z\) is a multiple of \(Q\), all remainder coordinates vanish, and
every integer matrix satisfies \(F(Uz)=[UF(z)]_K\). This does not
contradict Theorem 1, whose hypothesis ranges over all inputs.

For an arbitrary subset \(\Omega\subseteq\mathbb Z_P^m\), a
factorization exists on \(\Omega\) exactly when \(F(Uz+b)\) is constant
on each nonempty set \(\Omega\cap F^{-1}(y)\). Necessity is immediate;
for sufficiency, define \(g(y)\) by that constant and choose arbitrary
values on unused fibers. This is a set-theoretic criterion and supplies
no efficient algorithm by itself.

\subsection{5. Optimal prediction of a missing
carry}\label{optimal-prediction-of-a-missing-carry}

Let \(Z\) have a specified finite distribution, and set \(Y=F(Z)\) and
\(T=F(UZ+b)\). A simulator receives only \(Y\) and uses fresh
randomness, independent of \(Z\) conditional on \(Y\). Its smallest
disagreement probability is

\[
\varepsilon_* =1-\sum_y\max_t\Pr(Y=y,T=t).
\tag{4}
\]

For a kernel \(K(y,t)\), success is \(\sum_{y,t}\Pr(Y=y,T=t)K(y,t)\).
Maximizing each row of this linear expression puts all its mass on a
most likely target value. This proves (4). The optimum is
information-theoretic; computing the required conditional probabilities
may be difficult.

Now let \(Z_1,Z_2\) be independent and uniform on \(\mathbb Z_P\), and
set \(Y=(F(Z_1),F(Z_2))\) and \(T=F(Z_1+Z_2)\). The remainders
\(R_1,R_2\) are independent uniform elements of \(\{0,\ldots,Q-1\}\),
independent of the quotients. The missing carry
\(C=\lfloor(R_1+R_2)/Q\rfloor\) is binary, with

\[
p_0=\Pr(C=0)=\frac{Q+1}{2Q},\qquad p_1=\Pr(C=1)=\frac{Q-1}{2Q}.
\tag{5}
\]

There are \(Q(Q+1)/2\) pairs whose sum is below \(Q\). The two possible
output values differ modulo \(K\ge2\). Hence

\[
\varepsilon_* =\frac{Q-1}{2Q}.
\tag{6}
\]

The same optimum holds for the entire shear output, since its second
quotient is already known.

\subsection{6. Public simulation and joint
simulation}\label{public-simulation-and-joint-simulation}

We use total variation in the convention
\(\Delta(\mu,\nu)=\tfrac12\sum_x|\mu(x)-\nu(x)|\). In the preceding
example, independently drawing a new carry with probabilities
\((p_0,p_1)\) reproduces the law of \((Y,T)\) exactly. It need not
reproduce the dependence of \(T\) on the original unrounded input.

Whenever \(T\) is deterministic given \(Z\) and \(Y\) is a function of
\(Z\), a fresh \(Y\)-only simulator satisfies

\[
\Delta\bigl(\mathcal L(Z,T),\mathcal L(Z,\widetilde T)\bigr)
=\Pr(T\ne\widetilde T).
\tag{7}
\]

To prove this, condition on each \(z\): the distance from the point mass
at \(T(z)\) to the simulated conditional distribution is
\(1-\Pr(\widetilde T=T(z)\mid z)\). Averaging proves the identity. For
independent posterior resampling of the carry, (7) gives

\[
2p_0p_1=\frac{Q^2-1}{2Q^2},
\tag{8}
\]

whereas the smallest joint distance is \(p_1\), obtained by always
predicting the more likely carry. Exactness of a public marginal and
optimality for a stronger joint view are therefore different objectives.

A binary-secret example makes the distinction especially transparent.
Let \(P=4\), \(Q=2\), take the public column \(A=(1,1)^\mathsf{T}\), a
uniform bit \(S\), and no noise. The observed vector \(Y=F(AS)\) is
always \((0,0)\). Under the shear \((z_1,z_2)\mapsto(z_1+z_2,z_2)\), the
target is \(T=(S,0)\). Every simulator using only \((A,Y)\) is
independent of \(S\), so its distance from the joint law of \((S,T)\) is
at least \(1/2\). Equality is attained by every simulator supported on
\((0,0)\) and \((1,0)\); an equal mixture also reproduces the public
target marginal exactly. No choice of fresh public randomness
reconstructs the missing correlation with \(S\).

These examples do not exclude more general reductions with additional
information or altered distributions. They identify a specific
information constraint: \(Z\to Y\to\widetilde T\) is a Markov chain,
while the real target can retain information about \(Z\) beyond \(Y\).

\subsection{7. Relation to earlier work and
scope}\label{relation-to-earlier-work-and-scope}

Partition-preserving transformation semigroups provide the general
language for condition (1) {[}6{]}. The normal form above computes the
intersection of affine maps with transformations preserving a fixed
product of interval cells. Congruence-preserving functions {[}7{]} are
defined using algebraic congruences. Our cells are not cosets of a
proper subgroup: a cell containing both \(0\) and \(1\) cannot be such a
coset. For example, multiplication by \(2\) modulo \(4\) preserves group
congruences, but separates \(0\) and \(1\) in the same rounding cell for
\(Q=2\). Even translation by \(1\) can split a cell. These differences
prevent an automatic transfer of classifications for subgroup partitions
to this setting.

The LWE and LWR literature {[}1--5{]} supplies the cryptographic
motivation, not a hardness consequence of Theorem 1. The theorem
concerns universal exact postprocessing at fixed parameters. It does not
cover unequal cells when \(Q\) does not divide \(P\), nearest-integer
conventions, changes of moduli or dimensions, adaptive side information,
or simulations required only on a restricted secret-dependent support.
The carry identity, finite prediction formula, and marginal-versus-joint
distinction are standard tools; the specific affine classification and
sharp examples constitute the contribution of this note.

\subsection{8. Reproducible finite
verification}\label{reproducible-finite-verification}

The supplementary program \texttt{scripts/verify.py} uses integer
arithmetic and \texttt{fractions.Fraction}. It checks all affine rows
and translations for proper divisors \(Q\ge2\) of \(4\le P\le12\), in
dimensions \(m=1,2\), and also \(m=3\) for \(P\le6\). Direct constancy
on the zero cell is compared with the stated normal form; the verifier
finds a collision witness for every excluded affine row map. This covers
\textbf{15,082 affine rows}, \textbf{1,923,906 evaluated points}, and
\textbf{14,578 excluded maps for which a collision is found}. The
predicted number of admissible affine rows for a fixed domain is
\(P+2mK\).

Further checks verify (2) in \textbf{1,451,375} signed arithmetic cases
and verify the prediction and joint-distance formulas for \textbf{45}
laws with \(2\le Q\le16\) and \(K\in\{2,3,4\}\). Reflection and
restricted-support examples are also checked. The binary-secret proof is
analytic; its additional probability-grid check is only a finite
corroboration. Complete outputs, parameter domains, dependencies, and
reproduction commands accompany the paper. These computations do not
replace the proofs or establish any cryptographic security level.

\subsection{References}\label{references}

{[}1{]} Oded Regev. \emph{The Learning with Errors Problem}. IEEE
Conference on Computational Complexity, 2010, pp.~191--204.
\href{https://cims.nyu.edu/~regev/papers/lwesurvey.pdf}{Author version}.

{[}2{]} Abhishek Banerjee, Chris Peikert, and Alon Rosen.
\emph{Pseudorandom Functions and Lattices}. EUROCRYPT 2012.
\href{https://web.eecs.umich.edu/~cpeikert/pubs/prf-lattice.pdf}{Author
version}.

{[}3{]} Joël Alwen, Stephan Krenn, Krzysztof Pietrzak, and Daniel Wichs.
\emph{Learning with Rounding, Revisited: New Reduction, Properties and
Applications}. CRYPTO 2013, pp.~57--74.
\href{https://crypto.ethz.ch/publications/files/AKPW13.pdf}{Author
version}. DOI: 10.1007/978-3-642-40041-4\_4.

{[}4{]} Andrej Bogdanov, Siyao Guo, Daniel Masny, Silas Richelson, and
Alon Rosen. \emph{On the Hardness of Learning with Rounding over Small
Modulus}. TCC 2016-A.
\href{https://people.csail.mit.edu/sirichel/lwr.pdf}{Author version}.

{[}5{]} Shi Bai, Katharina Boudgoust, Dipayan Das, Adeline
Roux-Langlois, Weiqiang Wen, and Zhenfei Zhang. \emph{Middle-Product
Learning with Rounding Problem and its Applications}. ASIACRYPT 2019,
pp.~55--81. \href{https://eprint.iacr.org/2019/1001}{IACR ePrint
2019/1001}.

{[}6{]} Mosarof Sarkar and Shubh N. Singh. \emph{On certain semigroups
of transformations that preserve a partition}. Communications in Algebra
49 (2021), 331--342.
\href{https://doi.org/10.1080/00927872.2020.1799224}{DOI:
10.1080/00927872.2020.1799224}.
\href{https://arxiv.org/abs/2006.04242}{Author version:
arXiv:2006.04242}.

{[}7{]} C. J. Maxson and Frederik Saxinger. \emph{Rings of congruence
preserving functions}. Monatshefte für Mathematik 187 (2018), 531--542.
\href{https://doi.org/10.1007/s00605-017-1105-3}{DOI:
10.1007/s00605-017-1105-3}.

{[}8{]} Yuanxin Guo, Stark C. Draper, and Wei Yu. \emph{Modulo
Quantization Coding for Gaussian Primitive Relay Channel with Perfectly
Correlated Noises}. IEEE International Symposium on Information Theory,
2025.
\href{https://www.comm.utoronto.ca/~weiyu/2025_ISIT_Yuanxin.pdf}{Author
version}.

\end{document}
